% - General formulation: for each class $y \in \mathcal{Y}$, define a transformation $\phi_y$:$[0,1]^K \to [0,1]^K$ that maps 
%   the raw score vector to an conformal score vector such that a true infector of class $y$ is near the origin\\
% - Property we require: $phi_y(\mathbf{s})$ is small when $\mathbf{s}$ is consistent with label $y$, large otherwise\\
% - we make a single unified envelope coherent across all classes

\subsection{Nonconformity Score Transformation}
\label{sec3.3: nc transformation}
\noindent The raw score vectors are not directly comparable across candidate labels and cannot be used directly as nonconformity 
measures. Recall from Section~\ref{sec2.4.1: ensemble score vectors} that the CSA framework requires a label-dependent 
transformation $\phi_h: ([0,1] \cup \{\texttt{NA}\})^K \times \{0,1\}^K \to \left(\mathbb{R}_+ \cup \{\texttt{NA}\}\right)^K$ 
that maps candidate hypotheses into a  nonconformity score space centered near the origin. 

\vspace{0.15cm}
\noindent The raw functional category scores $s_k \in [0, 1]$ quantify evidence in favor of functional compatibility, therefore a raw 
score close to $1$ reflects strong infection compatibility (low nonconformity), whereas a score near $0$ reflects high 
nonconformity. Thus, the mapping is defined such that when a candidate hypothesis $h \in \mathcal{H}$ corresponds to the true 
host label, the resulting vector $\mathbf{v} = \phi_h(\mathbf{s}, \mathbf{m})$ concentrates near $\mathbf{0} 
\in \mathbb{R}_+^K$, whereas incorrect hypotheses yield vectors situated significantly further from the origin. We now 
formalize $\phi_h$ for both classification settings while explicitly preserving structural missingness via the 
mask vector $\mathbf{m}$.

\vspace{0.15cm}
\noindent To map all the score vectors into one unified nonconformity space near the origin, we define label-dependent 
transformations $\phi_h$. Formally, for a single score dimension $s_k \in [0,1]$ with the binary mask $m_k \in \{0,1\}$, 
we define the general coordinate wise mapping:
\begin{equation}
\label{eq: general nc transform}
\phi_h(s_k, m_k) = 
\begin{cases} 
|s_k - c_h| & \text{if } m_k = 1, \\
\texttt{NA} & \text{if } m_k = 0,
\end{cases}
\qquad c_h \in \{0, 1\},
\end{equation}

\noindent where the $c_h$ tells us whether a low ($c_h = 0$) or high ($c_h = 1$) raw score provides the evidence for hypothesis $h$. 
$s_k \in [0,1]$, \eqref{eq: general nc transform} can be simplified to $\phi_h(s_k) = s_k$ when $c_h = 0$, and 
$\phi_h(s_k) = 1 - s_k$ when $c_h = 1$, this ensures that $\phi_h(s_k) \to 0$ whenever $s_k \to c_h$.


\subsubsection{Gram Type Classification Nonconformity Mapping}
\label{sec3.3.1: gram nc mapping}
\noindent In the binary Gram-type classification, $\mathcal{H} = \{\text{gram-neg}, \text{gram-pos}\}$, all the $K$ functional 
score models measure a single Gram-positive evidence scale. Empirically, true Gram-negative phages have low raw scores 
(mean $\approx 0.07$), whereas true Gram-positive phages exhibit high raw scores (mean $\approx 0.92$).

\noindent Since the scale is shared, the $c_h$ must differ by hypothesis: $c_{\text{gramneg}} = 0$ and $c_{\text{grampos}} = 1$. 
This gives us the elementwise transformation for dimension $k \in \{1, \ldots, K\}$:
\begin{equation}
\label{eq: gram nc transformation}
v_k = (\phi_h(\mathbf{s}, \mathbf{m}))_k = \begin{cases} 
s_k & \text{if } h = \text{gram-neg} \text{ and } m_k = 1, \\
1 - s_k & \text{if } h = \text{gram-pos} \text{ and } m_k = 1, \\
\texttt{NA} & \text{if } m_k = 0.
\end{cases}
\end{equation}

\noindent Using \eqref{eq: gram nc transformation}, we get a high nonconformity score if we test an incorrect hypothesis, and 
the scores for the true hypothesis are mapped near the origin $\mathbf{0} \in \mathbb{R}_+^K$. This allows both classes 
to be calibrated against a shared reference envelope $E \subset \mathbb{R}_+^K$.

\subsubsection{Host Level Classification Nonconformity Mapping}
\label{sec3.3.2: host nc mapping}
\noindent In the host level setting ($N > 2$), the candidate labels do not share the scores. Instead, each candidate host $h 
\in \mathcal{H}$ is evaluated on a host-specific sub-vector $\mathbf{s}_h \in ([0,1] \cup \{\text{NA}\})^{K_h}$. This measures
the direct functional compatibility with host $h$. 

\vspace{0.15cm}
\noindent A high raw score $s_{h,k}$ indicates high likelihood of infection, $c_h = 1$ for all hosts, using this we can get the 
following host-specific mapping across $k \in \{1, \ldots, K_h\}$:
\begin{equation}
\label{eq: host nc transformation}
v_{h,k} = (\phi_h(\mathbf{s}_h, \mathbf{m}_h))_k = \begin{cases} 
1 - s_{h,k} & \text{if } m_{h,k} = 1, \\
\text{NA} & \text{if } m_{h,k} = 0.
\end{cases}
\end{equation}

\noindent Unlike the Gram classification task, in host level prediction we construct $N$ separate host envelopes $E_h \subset 
\mathbb{R}_+^{K_h}$ which are fitted independently on calibration phages with true label $H = h$. At test time, $\mathbf{s}_h$ 
is transformed using \eqref{eq: host nc transformation} and tested against $E_h$, this avoids forcing unrelated host signatures 
into a single shared coordinate system.

\vspace{0.15cm}
\noindent Both transforms are strictly monotone fixed maps that require no parameter estimation from data. Therefore, they do 
not have sample dependencies that could break the exchangeability argument (which is necessary for the coverage guarantee, Theorem
~\ref{thm: marginal coverage}). The missingness masks $\mathbf{m}$ propagate unchanged as $\texttt{NA}$ values, and the 
envelope constructions of Section~\ref{sec3.4: envelope construction} operate directly on this representation.

